\section{Introduction}
Electronic-support and electronic-attack systems interact through a resource-constrained loop: an interferer chooses where and when to spend energy, while a multifunction radar chooses which service and bearing to observe before each mission deadline. A high instantaneous JNR can be operationally irrelevant if the radar detects the corresponding target in another dwell, whereas a modest interference pulse can matter when it arrives at the only vulnerable point in a deadline window. The operational outcome is therefore mission completion, not a single link-budget scalar.

This distinction exposes a gap between two established lines of work. Radar--jammer learning studies formulate waveform, frequency, or power decisions as sequential games. Multifunction-radar scheduling studies optimize scan and track allocation under time or power constraints. These lines of work leave open a system-level question: how does the balance of an adaptive game change when the number of coordinated attackers increases while the defender's mission workload and total interference budget remain fixed?

We answer this question with a controlled S6-to-S7 progression in FluxPhased. S6 contains one learned jammer and two learned radar heads. S7 contains two parameter-shared jammer agents and two parameter-shared radar agents. The S7 team receives a 63-cell activation budget, split 32/31 between jammers. This controls total jammer activation while the radar-site geometry changes from S6 to S7; the matched co-located-jammer control isolates the additional cross-fire component within the S7 geometry.

The main result is a change in containment, not a corresponding failure of learning. Across three S7 seeds trained to 3000 iterations, head-to-head mission drop is $0.3366\pm0.0143$, compared with $0.0860\pm0.0061$ in the valid three-seed 12-dB S6 baseline. The floor-adjusted fraction of marginal jammer disruption removed by learned radars falls from $62.8\%\pm1.6\%$ to $23.0\%\pm1.1\%$. The radar floor against idle jammers remains low ($0.0194\pm0.0008$ in S7 versus $0.0242\pm0.0097$ in valid S6), and the S7 trajectory is stationary from approximately iteration 1700 through 3000. Thus, the second attacker changes the stable equilibrium's scale while preserving low idle-jammer loss and training stability.

Evaluation-only scripted baselines bound these equilibrium numbers. Both learned teams clearly beat random and fixed-stare scripted opponents, but a greedy mission-stare radar holds mission drop at $0.0889$ against every trained jammer team---below the learned head-to-head value and below the natural scheduling floor. The head-to-head drop is therefore an equilibrium property of the co-adapted pair, not a ceiling on radar performance, and both teams are specialized to their training opponents.

This paper makes five contributions:
\begin{enumerate}
    \item We formulate a test-and-evaluation protocol for learned radar resource management under adaptive electronic attack, coupling deadline-constrained mission loss to array-factor link physics, finite energy, and two-team multi-agent PPO. Its three components---the pre-training contestability gate, the floor-adjusted containment metric $\eta$, and the evaluation-only scripted anchors---are each independently reusable test instruments, separable from the specific learner used here.
    \item We quantify a benchmark-level attacker-count scaling effect: with total jammer activation controlled, doubling attackers reduces the mission-denial containment of learned radar scheduling by roughly two-thirds, with replication across three converged S7 seeds, a three-seed valid S6 reference, and a three-seed co-located control.
    \item We use a matched co-located-jammer control to decompose the mechanism under the separated S7 radar geometry. Two jammers already reduce containment to $26.6\%\pm2.9\%$ across three control seeds; cross-fire adds a secondary reduction to $24.2\%$.
    \item We anchor both learned teams with evaluation-only scripted baselines. Both teams dominate random and stare policies, but a greedy mission-stare radar is unexploited by every trained jammer team, showing that head-to-head drop is an equilibrium quantity rather than a radar performance ceiling.
    \item We test a minimal opponent-class mixture. The singleton-to-pair disruption ratio falls from $0.502$ to $0.262$ as singleton exposure increases, identifying opponent-class mixing as a practical robustness intervention while documenting its training-budget confound.
\end{enumerate}

The rest of the paper defines the benchmark and evaluation metric, presents convergence and scaling results, tests the count--geometry decomposition, analyzes opponent-specific adaptation, and evaluates the lightweight mixing intervention. We use ``drop'' for the fraction of eligible missions that are not completed before their deadlines; higher drop therefore indicates greater jammer-side success.
